\chapter{考察}\label{cha:Evaluation}

\section{関連研究}
\section{評価実験}
\section{評価実験}

本研究で提案するプロンプトの有効性を検証するため、被験者を用いた比較実験を行った。
実験では、実際のシステム開発を想定した題材を用い、提案手法とその他の手法で作成された章立ての品質を専門家が評価した。

\subsection{実験設定}

\subsubsection{対象システム}
実験の題材として、以下の2つのシステムを選定した。
\begin{enumerate}
    \item \textbf{BMP画像分割アプリ}: デスクトップアプリケーションを想定（小規模・単機能）
    \item \textbf{クラウド型電子契約サービス}: Webアプリケーションを想定（中規模・業務機能多数）
\end{enumerate}

\subsubsection{比較手法}
作成手法として、以下の3パターンを設定した。
\begin{itemize}
    \item \textbf{手法A（手動作成）}: LLMを用いず、被験者の知識のみで章立てを作成する。
    \item \textbf{手法B（LLMのみ）}: LLMを使用するが、本研究のプロンプトは使用せず、自由な指示で章立てを作成する。
    \item \textbf{手法C（提案手法）}: LLMを使用し、本研究で作成したテンプレートプロンプトを用いて章立てを作成する。
\end{itemize}
なお、被験者CとDの手法Bにおいて、プロンプトの指定はないものとする。

\subsubsection{被験者と実験の割り当て}
4名の被験者（A, B, C, D）に対し、システムと手法の組み合わせを以下の表\ref{tb:experiment_allocation}のように割り当てた。
クロス集計的な評価が可能になるように、被験者ごとに異なる組み合わせを設定している。

\begin{table}[h]
  \caption{被験者ごとのシステムと採用手法の割り当て}
  \label{tb:experiment_allocation}
  \centering
  \begin{tabular}{|c|l|l|}
    \hline
    \textbf{被験者} & \textbf{BMP画像分割アプリ} & \textbf{クラウド型電子契約サービス} \\
    \hline\hline
    A & 手法A（手動） & 手法C（提案手法） \\
    \hline
    B & 手法C（提案手法） & 手法A（手動） \\
    \hline
    C & 手法B（LLMのみ） & 手法C（提案手法） \\
    \hline
    D & 手法C（提案手法） & 手法B（LLMのみ） \\
    \hline
  \end{tabular}
\end{table}

\subsection{評価基準}
作成された章立ての成果物は、株式会社スカイコムにて5年以上のシステム仕様書作成経験を持つ専門家によって評価された。
評価は以下の8項目について行われ、それぞれ5段階のリッカート尺度（5: 非常に良い 〜 1: 非常に悪い）で採点された。

\begin{enumerate}
    \item 最低限含むべき章が含まれているか（網羅性）
    \item 章の順序は適切か（論理性）
    \item 章タイトルは適切か（命名の適切さ）
    \item 章立て全体として理解しやすいか（可読性）
    \item 章の粒度は適切か（詳細度）
    \item 想定読者（新人プログラマ等）に適した章立てか（適合性）
    \item 将来的な拡張に対応しやすい章立てか（拡張性）
    \item 章、節、項の包含関係は適切か（構造的整合性）
\end{enumerate}

\subsection{実験結果}
% ここに結果を記述
\section{考察}
\section{課題}

\iffalse
本章では、論文執筆における章立てとセクション構成について説明する。

% \section{章の構成}

論文は複数の章で構成される。各章は \verb|\chapter{}| コマンドで定義する。

% \subsection{章番号の自動付与}

LaTeX では章番号が自動的に付与される。
例えば、本章は第\ref{cha:Evaluation}章である。

% \subsection{章のラベル付け}

各章には \verb|\label{}| を付けることで、他の箇所から参照できる。
これにより、章番号が変更されても参照が自動的に更新される。

% \section{セクションの構成}

各章は複数のセクションで構成される。

% \subsection{セクションの階層}

LaTeX では以下の階層構造が使用できる：

\begin{enumerate}
  \item \verb|\chapter{}|: 章 (1, 2, 3, ...)
  \item \verb|\section{}|: 節 (1.1, 1.2, 1.3, ...)
  % \item \verb|\subsection{}|: 小節 (1.1.1, 1.1.2, 1.1.3, ...)
  \item \verb|\subsubsection{}|: 小小節 (1.1.1.1, 1.1.1.2, 1.1.1.3, ...)
\end{enumerate}

% \subsection{適切な階層の選択}

論文の構成を考える際は、以下の点に注意する：

\begin{itemize}
  \item 階層が深くなりすぎないようにする（通常は subsection まで）
  \item 各セクションには適切な量の内容を含める
  \item セクションの粒度を統一する
\end{itemize}

% \section{論文の一般的な構成}

片山徹郎研究室で主に使われてきた、修士論文の一般的な構成は以下の通りである：

\begin{enumerate}
  \item はじめに
  \item 研究の準備（使用ツール、ライブラリの紹介）
  \item 機能
  \item 実装
  \item 適用例
  \item 考察
  \item おわりに（まとめと今後の課題）
\end{enumerate}

この構成が全てではなく、研究内容に合わせて適宜変更されることが多い。
例えば、卒業論文で開発したツールの拡張を行う場合は、拡張した箇所の説明を追加することが多い。
詳細な構成などは、片山研の \verb|graduate| フォルダにある論文を参考にすると良い。

% \section{本テンプレートの構成}

本テンプレートでは、以下の構成でサンプルを提供している：

\begin{itemize}
  \item 第\ref{cha:Introduction}章: テンプレートの使い方
  \item 第\ref{cha:Preparation}章: 基本的な LaTeX の書き方
  \item 第\ref{cha:Function}章: 図、表、数式の挿入方法
  \item 第\ref{cha:Implementation}章: ソースコードの挿入方法
  \item 第\ref{cha:Indication}章: 参考文献の引用方法
  \item 第\ref{cha:Evaluation}章: 章立てとセクション構成
  \item 第\ref{cha:Conclusion}章: まとめ
\end{itemize}

実際の論文執筆時は、各章の内容を自分の研究内容に置き換えて使用する。
\fi